% ============================================================
%  CHAPITRE 12 — Tests et validation
%  À inclure dans le document principal avec :
%    \input{chapitre12}   ou   \include{chapitre12}
% ============================================================

\chapter{Tests et validation}
\label{chap:tests-validation}

La qualité d'un système logiciel ne se décrète pas~: elle se démontre. Ce chapitre
présente l'ensemble de la démarche de validation mise en œuvre pour BCaaS, depuis la
stratégie globale de test jusqu'au scénario de validation bout-en-bout, en passant par
les critères d'acceptation formels. L'objectif est de montrer que chaque brique du système
--- sécurité, logique métier, isolation multi-tenant, génération de factures --- a été
vérifiée de manière systématique et reproductible.

% ------------------------------------------------------------
\section{Stratégie de test}
\label{sec:strategie-test}

% ---- 12.1.1 ------------------------------------------------
\subsection{Vue d'ensemble et outillage}
\label{subsec:outillage}

La validation de BCaaS repose sur \textbf{16~classes de test} totalisant
\textbf{152~méthodes}, organisées selon une pyramide de test à trois niveaux. L'ensemble
du dispositif s'appuie sur la stack de test standard de l'écosystème Spring~Boot~:

\begin{description}[leftmargin=0.5cm]
  \item[\textbf{JUnit~5} (\textit{Jupiter})]
    Socle d'exécution des tests~: annotations \texttt{@Test}, \texttt{@BeforeEach},
    \texttt{@AfterEach}, \texttt{@ParameterizedTest}. JUnit~5 apporte la prise en charge
    des extensions (\texttt{@ExtendWith}), utilisée notamment pour injecter le contexte
    Spring ou Mockito.

  \item[\textbf{Mockito}]
    Bibliothèque de création de \textit{mocks} et de \textit{stubs}. Elle permet de
    simuler les dépendances externes d'un composant (repositories, services tiers) afin de
    tester la logique d'un service en totale isolation, sans base de données réelle ni
    appel réseau.

  \item[\textbf{AssertJ}]
    Bibliothèque d'assertions fluides (\texttt{assertThat(...).isEqualTo(...)},
    \texttt{.contains(...)}, \texttt{.isNotNull()}...). Sa syntaxe chaînée améliore la
    lisibilité des assertions et produit des messages d'erreur plus explicites que les
    assertions JUnit natives.

  \item[\textbf{spring-boot-starter-test}]
    Méta-dépendance qui embarque JUnit~5, Mockito et AssertJ, ainsi que
    \texttt{MockMvc} pour les tests de contrôleurs HTTP et \texttt{@SpringBootTest} pour
    les tests d'intégration avec contexte Spring complet.

  \item[\textbf{spring-security-test}]
    Extension dédiée aux tests de la couche de sécurité Spring Security. Elle fournit les
    annotations \texttt{@WithMockUser} et \texttt{@WithUserDetails}, ainsi que les
    utilitaires \texttt{SecurityMockMvcRequestPostProcessors} pour simuler des requêtes
    authentifiées sans générer de vrais tokens JWT.
\end{description}

% ---- 12.1.2 ------------------------------------------------
\subsection{Environnement de test}
\label{subsec:env-test}

Les tests s'exécutent dans un environnement entièrement isolé de la base de données de
production ou de développement. Deux choix techniques fondamentaux garantissent cette
isolation~:

\paragraph{Base H2 en mémoire.}
H2 est un moteur de base de données relationnelle embarqué, écrit en Java, qui s'exécute
entièrement en mémoire vive (\textit{in-memory}). Il n'est pas nécessaire d'installer ni
de démarrer un serveur de base de données externe~: H2 démarre automatiquement avec le
contexte Spring et s'arrête à la fin de la suite de tests. La configuration
\texttt{ddl-auto=create-drop} garantit que le schéma est recréé à chaque lancement et
détruit à la fin, rendant chaque exécution de tests parfaitement reproductible et
indépendante de l'état laissé par une exécution précédente.

\paragraph{Liquibase désactivé.}
En production et en développement, BCaaS utilise Liquibase pour gérer les migrations de
schéma de façon versionnée. Liquibase est désactivé dans le profil de test~: le schéma est
entièrement reconstruit par Hibernate à partir des entités JPA (\texttt{create-drop}), ce
qui garantit que les tests ne dépendent pas d'un état de migration particulier et
s'exécutent aussi bien sur la première version que sur une version ultérieure du schéma.

\paragraph{Profil Spring dédié.}
Un profil Spring \texttt{test} est activé pour les tests d'intégration. Ce profil
surcharg les propriétés de connexion (\texttt{spring.datasource.url=jdbc:h2:mem:testdb}),
désactive les mécanismes de cache, et configure des secrets JWT simplifiés adaptés à
l'environnement de test.

% ---- 12.1.3 ------------------------------------------------
\subsection{Les trois niveaux de test}
\label{subsec:niveaux-test}

La stratégie de test de BCaaS suit le principe de la \textbf{pyramide de test}~: une large
base de tests unitaires rapides, une couche intermédiaire de tests de contrôleurs, et un
sommet de tests d'intégration plus lents mais plus représentatifs du comportement réel du
système.

\subsubsection*{Niveau 1 — Tests unitaires de services}

Les tests unitaires ciblent la \textbf{logique métier pure} des services applicatifs, en
isolation complète de toute infrastructure. Toutes les dépendances (repositories,
services tiers) sont remplacées par des mocks Mockito.

Les services couverts sont~:

\begin{itemize}
  \item \textbf{InvoiceService}~: génération automatique de facture lors de la transition
        \texttt{fulfill}, calcul des montants, association à la demande de service et au
        tenant.
  \item \textbf{BusinessService}~: création, modification et suppression de métiers~;
        vérification des règles de nommage et d'unicité au sein d'un tenant.
  \item \textbf{TenantService}~: provisionnement d'un nouveau tenant, vérification de
        l'absence de doublon, initialisation des configurations par défaut.
  \item \textbf{ActorService}~: gestion du cycle de vie des acteurs (prestataires, clients)
        et de leurs rôles au sein d'un tenant.
  \item \textbf{ServiceRequestService}~: machine à états de la demande de service~;
        validation des transitions autorisées et déclenchement des effets de bord
        (notification, facturation).
  \item \textbf{ServiceCatalogueService}~: gestion des services offerts par un métier~;
        règles de visibilité et de tarification.
  \item \textbf{BusinessRuleService}~: évaluation des règles métier configurables par
        tenant~; vérification de la conformité des demandes aux politiques en vigueur.
  \item \textbf{TenantContextTest}~: vérification du bon fonctionnement du mécanisme de
        propagation du \texttt{tenant\_id} dans le contexte de thread (ThreadLocal)~; test
        des cas de corruption ou d'absence de contexte.
\end{itemize}

\subsubsection*{Niveau 2 — Tests de contrôleurs (WebMvc)}

Les tests de contrôleurs utilisent \texttt{MockMvc}, qui simule des requêtes HTTP sans
démarrer un vrai serveur web. Le contexte Spring est chargé partiellement
(\texttt{@WebMvcTest}), incluant uniquement la couche web (contrôleurs, filtres,
sérialiseurs JSON) et excluant la couche de persistance.

Les services sont mockés via \texttt{@MockBean}, ce qui permet de se concentrer sur~:

\begin{itemize}
  \item la \textbf{désérialisation} des corps de requêtes JSON et la validation des
        contraintes (\texttt{@Valid}, \texttt{@NotNull}, tailles, formats)~;
  \item le \textbf{routage} des URLs vers les bonnes méthodes de contrôleur~;
  \item la \textbf{sérialisation} des réponses JSON et les codes HTTP retournés (200, 201,
        400, 403, 404, 409)~;
  \item la \textbf{gestion des erreurs} via les \texttt{@ExceptionHandler} et le
        \texttt{ControllerAdvice} global.
\end{itemize}

Les contrôleurs couverts sont \textbf{BusinessController}, \textbf{InvoiceController},
\textbf{ActorController} et \textbf{ServiceRequestController}.

\subsubsection*{Niveau 3 — Tests de sécurité et d'intégration}

Ce niveau constitue le sommet de la pyramide~: les tests chargent le contexte Spring
complet (\texttt{@SpringBootTest}) avec la base H2 en mémoire, et exercent le système
de bout en bout à travers plusieurs couches.

\begin{description}[leftmargin=0.5cm]
  \item[\textbf{JwtServiceTest}]
    Vérifie la génération, la signature et la validation des tokens JWT~: format du token,
    présence et valeur des \textit{claims} (\texttt{sub}, \texttt{tenantId}, \texttt{exp},
    \texttt{iat}), rejet des tokens expirés ou falsifiés, et extraction correcte du
    \texttt{tenantId} depuis le token.

  \item[\textbf{SecurityIntegrationTest}]
    Teste le filtre de sécurité sur l'ensemble des endpoints protégés~: rejet des requêtes
    sans token (401), rejet des tokens invalides (401), autorisation correcte des requêtes
    avec un token valide, et vérification que les endpoints publics (inscription, connexion)
    sont bien accessibles sans authentification.

  \item[\textbf{ServiceRequestIntegrationTest}]
    Scénario d'intégration complet de la demande de service~: création d'un contexte tenant,
    d'un acteur demandeur et d'un prestataire, d'un service au catalogue, émission d'une
    demande, puis parcours de toutes les transitions de la machine à états
    (\texttt{pending} → \texttt{accepted} → \texttt{in\_progress} → \texttt{fulfilled})
    avec vérification, à chaque étape, de l'état persisté en base et de la génération
    automatique de la facture lors du passage à \texttt{fulfilled}.
\end{description}

% ------------------------------------------------------------
\section{Critères d'acceptation}
\label{sec:criteres-acceptation}

% ---- 12.2.1 ------------------------------------------------
\subsection{Définition et méthode}
\label{subsec:criteres-def}

Les critères d'acceptation formalisent les conditions minimales que BCaaS doit satisfaire
pour être considéré comme livrable. Chaque critère est identifié par un code
\texttt{ACC-XXX}, formulé de manière vérifiable et associé à une méthode de vérification
concrète. Cette approche s'inspire des pratiques de \textit{Behavior-Driven Development}
(BDD) et garantit que la validation n'est pas subjective~: un critère est soit satisfait,
soit non satisfait, sans ambiguïté.

% ---- 12.2.2 ------------------------------------------------
\subsection{Tableau des critères}
\label{subsec:criteres-tableau}

\begin{table}[ht]
  \centering
  \caption{Critères d'acceptation et méthodes de vérification}
  \label{tab:criteres-acceptation}
  \renewcommand{\arraystretch}{1.4}
  \begin{tabularx}{\textwidth}{>{\bfseries}p{1.8cm}
                               >{\raggedright\arraybackslash}X
                               >{\raggedright\arraybackslash}p{4.2cm}}
    \toprule
    \textbf{ID} & \textbf{Critère} & \textbf{Méthode de vérification} \\
    \midrule
    ACC-001 & Les endpoints REST répondent correctement aux requêtes valides et retournent
              les codes HTTP appropriés (2xx, 4xx) pour les cas d'erreur.
            & Tests WebMvc + validation manuelle via Swagger UI \\
    ACC-002 & L'authentification par JWT fonctionne~: génération, signature, validation et
              rejet des tokens expirés ou malformés.
            & \texttt{JwtServiceTest}, \texttt{SecurityIntegrationTest} \\
    ACC-003 & Le cycle complet CdS $\rightarrow$ FdS (demande de service à facture) se
              déroule sans erreur à travers toutes les transitions de la machine à états.
            & \texttt{ServiceRequestIntegrationTest} \\
    ACC-004 & La facture est générée automatiquement lors du passage à l'état
              \texttt{fulfilled}, avec les montants et références corrects.
            & Test de service \texttt{InvoiceServiceTest} + vérification en base H2 \\
    ACC-005 & L'isolation multi-tenant est garantie~: un tenant ne peut accéder aux données
              d'un autre, même en cas de manipulation du token.
            & \texttt{TenantContextTest} + \texttt{SecurityIntegrationTest} \\
    ACC-006 & La documentation Swagger (OpenAPI) est à jour et reflète fidèlement tous les
              endpoints exposés.
            & Validation visuelle sur \texttt{/swagger-ui.html} \\
    \bottomrule
  \end{tabularx}
\end{table}

% ---- 12.2.3 ------------------------------------------------
\subsection{Analyse des critères}
\label{subsec:criteres-analyse}

\paragraph{ACC-001 — Endpoints REST.}
Ce critère couvre la surface d'API complète de BCaaS. Les tests WebMvc vérifient
systématiquement les cas nominaux (requête valide $\rightarrow$ réponse 200/201) et les
cas d'erreur les plus fréquents~: corps de requête invalide (400), ressource inexistante
(404), conflit d'identifiant (409), et accès non autorisé (403). La validation manuelle
via Swagger UI complète ces tests automatisés en permettant une exploration visuelle
des contrats d'API.

\paragraph{ACC-002 — Authentification JWT.}
La sécurité de BCaaS repose entièrement sur les JWT~: chaque requête vers un endpoint
protégé doit porter un token valide, signé avec la clé secrète du serveur, non expiré, et
contenant un \texttt{tenantId} cohérent avec les données demandées. Les tests de ce critère
couvrent le chemin nominale (token valide $\rightarrow$ accès accordé) et les cas
d'attaque les plus courants (token absent, expiré, signature altérée, \texttt{tenantId}
modifié).

\paragraph{ACC-003 — Cycle CdS $\rightarrow$ FdS.}
Ce critère valide le flux métier central de BCaaS. CdS désigne la \textit{Commande de
Service} (demande émise par un client) et FdS la \textit{Facture de Service} générée à
l'issue de la prestation. Le test d'intégration associé parcourt toutes les transitions de
la machine à états et vérifie à chaque étape que~: (i) la transition est acceptée ou
refusée conformément aux règles métier, (ii) l'état persisté en base est correct, et (iii)
les effets de bord attendus (notifications, génération de facture) sont bien déclenchés.

\paragraph{ACC-004 — Génération automatique de la facture.}
La facture est le document qui matérialise la fin d'une prestation et déclenche le
processus de paiement. Sa génération automatique à l'état \texttt{fulfilled} est un
invariant métier critique~: aucune prestation ne doit pouvoir être marquée comme terminée
sans qu'une facture valide soit créée. Les tests vérifient également que les montants
calculés correspondent aux tarifs du catalogue et que les références croisées
(identifiant de la demande, identifiant du tenant, identifiant du prestataire) sont
correctement renseignées.

\paragraph{ACC-005 — Isolation multi-tenant.}
L'isolation multi-tenant est à la fois un critère de sécurité et un critère de conformité.
\texttt{TenantContextTest} vérifie que le mécanisme de propagation du \texttt{tenant\_id}
par ThreadLocal fonctionne correctement dans tous les cas~: contexte présent, contexte
absent, changement de contexte entre deux requêtes. \texttt{SecurityIntegrationTest}
complète cette vérification en s'assurant qu'un token portant un \texttt{tenantId} A ne
permet pas d'accéder aux ressources d'un tenant B.

\paragraph{ACC-006 — Documentation Swagger.}
La documentation Swagger (générée via SpringDoc/OpenAPI~3) constitue le contrat d'API
exposé aux développeurs consommateurs de BCaaS. Sa validité est vérifiée manuellement~:
chaque endpoint documenté est testé via l'interface Swagger UI pour s'assurer que les
paramètres, les corps de requête et les réponses décrits correspondent au comportement
réel de l'API.

% ------------------------------------------------------------
\section{Scénario de validation bout-en-bout}
\label{sec:scenario-e2e}

% ---- 12.3.1 ------------------------------------------------
\subsection{Objectif et portée}
\label{subsec:e2e-objectif}

Le scénario de validation bout-en-bout constitue la démonstration de référence de BCaaS.
Son objectif est de valider, en une seule séquence cohérente, l'ensemble des mécanismes
fondamentaux du système~: authentification JWT, isolation multi-tenant, machine à états
de la demande de service, et génération automatique de facture. Ce scénario est conçu
pour être reproductible aussi bien par les tests automatisés (\texttt{ServiceRequestIntegrationTest})
que par une démonstration manuelle lors d'une soutenance.

% ---- 12.3.2 ------------------------------------------------
\subsection{Déroulement du scénario nominal}
\label{subsec:e2e-scenario}

Le scénario se déroule en six étapes séquentielles, chacune produisant un état observable
vérifiable~:

\begin{enumerate}[leftmargin=*, label=\textbf{Étape \arabic*.}]

  \item \textbf{Inscription et connexion — obtention du JWT.}\\
    Un nouvel utilisateur s'inscrit via \texttt{POST /api/auth/register} en fournissant
    ses informations et le nom de son organisation (tenant). BCaaS provisionne
    automatiquement le tenant, crée le compte utilisateur et retourne un token JWT signé.
    Ce token contient le \texttt{tenantId} du tenant nouvellement créé et est inclus dans
    l'en-tête \texttt{Authorization: Bearer <token>} de toutes les requêtes suivantes.\\
    \textit{Vérification~:} le token est bien formé (trois segments base64 séparés par des
    points), non expiré, et le \texttt{tenantId} est présent dans les \textit{claims}.

  \item \textbf{Création d'un acteur, d'un portfolio et d'un métier.}\\
    L'utilisateur crée un acteur prestataire via \texttt{POST /api/actors}, lui associe
    un portfolio de compétences, puis crée un métier (\textit{business}) via
    \texttt{POST /api/businesses}. Ces entités sont automatiquement associées au tenant
    extrait du JWT.\\
    \textit{Vérification~:} les réponses HTTP retournent le statut 201, les entités
    créées sont récupérables via les endpoints \texttt{GET} correspondants et portent bien
    le \texttt{tenantId} attendu.

  \item \textbf{Ajout d'un service au catalogue.}\\
    Un service est ajouté au catalogue du métier via
    \texttt{POST /api/service-catalogue}. La définition du service inclut
    son libellé, sa description, son tarif unitaire et son unité de facturation.\\
    \textit{Vérification~:} le service est visible dans le catalogue du tenant via
    \texttt{GET /api/service-catalogue} et son tarif est correctement persisté.

  \item \textbf{Émission d'une demande de service (CdS).}\\
    Un acteur client émet une demande de service via
    \texttt{POST /api/service-requests}, en référençant le service du catalogue et en
    précisant les paramètres de la prestation (quantité, date souhaitée, instructions
    spécifiques). La demande est créée à l'état \texttt{PENDING}.\\
    \textit{Vérification~:} la demande est persistée avec l'état \texttt{PENDING},
    l'identifiant de la demande est retourné dans la réponse, et la demande est visible
    dans la liste des demandes du tenant.

  \item \textbf{Transitions d'état et génération de la facture.}\\
    Le prestataire fait progresser la demande à travers les transitions de la machine à
    états~:
    \begin{itemize}
      \item \texttt{PENDING} $\rightarrow$ \texttt{ACCEPTED}~: le prestataire accepte la
            demande via \texttt{PATCH /api/service-requests/\{id\}/accept}.
      \item \texttt{ACCEPTED} $\rightarrow$ \texttt{IN\_PROGRESS}~: le prestataire démarre
            la prestation via \texttt{PATCH /api/service-requests/\{id\}/start}.
      \item \texttt{IN\_PROGRESS} $\rightarrow$ \texttt{FULFILLED}~: le prestataire
            marque la prestation comme terminée via
            \texttt{PATCH /api/service-requests/\{id\}/fulfill}.
    \end{itemize}
    Lors du passage à \texttt{FULFILLED}, BCaaS génère automatiquement une facture
    (FdS) associée à la demande. Cette génération constitue l'\textbf{accusé de réception
    métier}~: elle matérialise la fin de la prestation et déclenche le processus de
    paiement.\\
    \textit{Vérification~:} après chaque transition, l'état persisté est contrôlé. Après
    \texttt{fulfill}, la facture est récupérable via \texttt{GET /api/invoices} et ses
    montants correspondent au tarif du catalogue multiplié par la quantité de la demande.

  \item \textbf{Paiement de la facture.}\\
    L'acteur client enregistre le paiement de la facture via
    \texttt{PATCH /api/invoices/\{id\}/pay}. La facture passe à l'état \texttt{PAID}
    et la demande de service est clôturée.\\
    \textit{Vérification~:} la facture est à l'état \texttt{PAID}, la date de paiement
    est renseignée, et aucune transition supplémentaire n'est acceptée par la machine à
    états (idempotence du paiement).

\end{enumerate}

% ---- 12.3.3 ------------------------------------------------
\subsection{Ce que le scénario valide}
\label{subsec:e2e-validation}

Ce parcours en six étapes exerce simultanément les mécanismes les plus critiques de
BCaaS~:

\begin{description}[leftmargin=0.5cm]
  \item[\textbf{La machine à états de la demande de service}]
    Chaque transition est validée~: les transitions autorisées aboutissent, les transitions
    interdites sont rejetées avec un code HTTP~409. L'ordre des transitions est imposé par
    la machine à états et ne peut pas être contourné via l'API.

  \item[\textbf{L'isolation multi-tenant}]
    Toutes les entités créées (acteurs, métiers, services, demandes, factures) sont
    automatiquement scopées au tenant du JWT. Un deuxième tenant créé en parallèle ne
    voit aucune de ces données.

  \item[\textbf{La génération automatique de facture}]
    La facture est créée de façon atomique avec la transition \texttt{fulfill}~: les deux
    opérations font partie de la même transaction JPA, garantissant qu'il ne peut pas
    exister de demande \texttt{FULFILLED} sans facture associée.

  \item[\textbf{La cohérence des références croisées}]
    La facture référence correctement la demande, le prestataire, le client et le tenant.
    Ces références sont vérifiées via des assertions sur les identifiants retournés par
    les différents endpoints.
\end{description}

% ------------------------------------------------------------
\section{Note sur la couverture de code et perspectives}
\label{sec:couverture}

% ---- 12.4.1 ------------------------------------------------
\subsection{État actuel}
\label{subsec:couverture-actuel}

Aucun outil de mesure de couverture de code (\textit{code coverage}) n'est actuellement
intégré au pipeline de build de BCaaS. Les 152~tests constituent néanmoins une base de
non-régression substantielle qui couvre~:

\begin{itemize}
  \item l'ensemble des chemins nominaux de tous les services métier~;
  \item les principaux cas d'erreur et cas limites identifiés durant le développement~;
  \item le cycle de vie complet de la demande de service, de l'émission au paiement~;
  \item les mécanismes de sécurité (authentification, autorisation, isolation tenant).
\end{itemize}

En l'absence de mesure formelle, il est impossible d'affirmer un taux de couverture précis.
Cependant, la couverture fonctionnelle --- c'est-à-dire la proportion des fonctionnalités
métier couvertes par au moins un test --- est estimée à un niveau satisfaisant pour un MVP.

% ---- 12.4.2 ------------------------------------------------
\subsection{Perspectives d'amélioration}
\label{subsec:couverture-perspectives}

Plusieurs améliorations sont identifiées pour les versions futures de BCaaS (voir
chapitre~14)~:

\begin{description}[leftmargin=0.5cm]
  \item[\textbf{Intégration de JaCoCo}]
    JaCoCo (\textit{Java Code Coverage Library}) s'intègre nativement avec Maven et
    Gradle. Son ajout permettrait de générer un rapport de couverture à chaque build,
    d'identifier les branches non testées et de définir des seuils minimaux de couverture
    (par exemple, 80\,\% de couverture de branches) en dessous desquels le build échoue.

  \item[\textbf{Tests de mutation}]
    Les tests de mutation (avec des outils comme PIT/Pitest) permettent d'évaluer la
    \textit{qualité} des tests plutôt que leur simple existence~: ils introduisent
    délibérément des fautes dans le code (mutations) et vérifient que les tests les
    détectent. Un test qui passe même en présence d'une mutation est un test insuffisant.

  \item[\textbf{Tests de charge}]
    BCaaS étant un service multi-tenant destiné à héberger plusieurs organisations
    simultanément, des tests de charge avec des outils comme Gatling ou JMeter permettraient
    de valider les performances sous des montées en charge réalistes et d'identifier les
    goulots d'étranglement avant la mise en production.

  \item[\textbf{Tests de contrat (Consumer-Driven Contract Testing)}]
    Si BCaaS est consommé par plusieurs applications clientes, les tests de contrat
    (avec Pact, par exemple) garantiraient que les évolutions de l'API ne cassent pas
    silencieusement les consommateurs existants.
\end{description}

% ------------------------------------------------------------
\section*{Synthèse}
\addcontentsline{toc}{section}{Synthèse}
\label{sec:synthese-chap12}

Ce chapitre a présenté la démarche complète de validation de BCaaS. La stratégie de test
repose sur une pyramide à trois niveaux --- tests unitaires, tests de contrôleurs et tests
d'intégration --- mobilisant 16~classes et 152~méthodes sur un environnement H2 en mémoire
parfaitement isolé. Les six critères d'acceptation formels garantissent que les
fonctionnalités les plus critiques --- sécurité JWT, isolation multi-tenant, cycle
CdS~$\rightarrow$~FdS --- sont vérifiées de manière systématique et reproductible.

Le scénario de validation bout-en-bout constitue la démonstration de référence du système~:
en six étapes cohérentes, il exerce simultanément la machine à états, l'isolation tenant
et la génération automatique de facture, offrant une preuve d'intégration concrète et
démontrable lors de la soutenance.

L'absence de mesure formelle de couverture constitue la principale limite reconnue de ce
dispositif~; son adressage, via l'intégration de JaCoCo et de tests de mutation, figure
parmi les priorités du plan d'évolution présenté au chapitre~14.